% Также можно использовать \Referat, как в оригинале
%\begin{abstract}
%	Титульный лист. Эта страница нужна мне, чтобы не сбивалась нумерация страниц
%	\cite{Dh}
%	\cite{Bayer}
%	\cite{Habr1}
%	\cite{Noise_func}
%	\cite{Ulich}

%Это пример каркаса расчётно-пояснительной записки, желательный к использованию в РПЗ проекта по курсу РСОИ.

%Данный опус, как и более новые версии этого документа, можно взять по адресу (\url{https://github.com/rominf/latex-g7-32}).

%Текст в документе носит совершенно абстрактный характер.
%\end{abstract}


\newenvironment{signstabular}[1][1]{
	\renewcommand*{\arraystretch}{#1}
	\tabular
}{
	\endtabular
}




\begin{titlepage}
	\thispagestyle{empty}
	
	\noindent\begin{minipage}{0.05\textwidth}
		\includegraphics[scale=0.3]{img/bmstu.png}
	\end{minipage}
	\hfill
	\begin{minipage}{0.85\textwidth}\raggedleft
		\begin{center}
			\fontsize{12pt}{0.3\baselineskip}\selectfont \textbf{Министерство науки и высшего образования Российской Федерации \\ Федеральное государственное бюджетное образовательное учреждение \\ высшего образования \\ <<Московский государственный технический университет \\ имени Н.Э. Баумана \\ (национальный исследовательский университет)>> \\ (МГТУ им. Н.Э. Баумана)}
		\end{center}
	\end{minipage}
	
	\begin{center}
		\fontsize{12pt}{0.1\baselineskip}\selectfont
		\noindent\makebox[\linewidth]{\rule{\textwidth}{4pt}} \makebox[\linewidth]{\rule{\textwidth}{1pt}}
	\end{center}
	
	\begin{flushleft}
		\fontsize{12pt}{0.8\baselineskip}\selectfont
		
		ФАКУЛЬТЕТ \uline{
			Информатика и системы управления
			\hfill}
		
		КАФЕДРА \uline{\mbox{\hspace{4mm}}
			Программное обеспечение ЭВМ и информационные технологии
			\hfill}
	\end{flushleft}
	
	\vfill
	
	\begin{center}
		\fontsize{20pt}{\baselineskip}\selectfont
		
		\textbf{РАСЧЕТНО-ПОЯСНИТЕЛЬНАЯ ЗАПИСКА}
		
		\textbf{\textit{К КУРСОВОМУ ПРОЕКТУ}}
		
		\textbf{\textit{НА ТЕМУ:}}
	\end{center}
	
	\begin{center}
		\fontsize{18pt}{0.6cm}\selectfont
		
		Разработка инструмента для построения трёхмерных изображений, ориентированного на микроконтроллеры STM32
		
	\end{center}
	
	\vfill
	
	\begin{table}[h!]
		\fontsize{12pt}{0.7\baselineskip}\selectfont
		\centering
		\begin{signstabular}[0.7]{p{7.25cm} >{\centering\arraybackslash}p{4cm} >{\centering\arraybackslash}p{4cm}}
			Студент группы \uline{ИУ7-53Б} & \uline{\mbox{\hspace*{4cm}}} & \uline{\hfill А.М. Сапожков  \hfill} \\
			& \scriptsize (Подпись, дата) & \scriptsize (И.О. Фамилия)
		\end{signstabular}
		
		\vspace{\baselineskip}
		
		\begin{signstabular}[0.7]{p{7.25cm} >{\centering\arraybackslash}p{4cm} >{\centering\arraybackslash}p{4cm}}
			Руководитель курсовой работы & \uline{\mbox{\hspace*{4cm}}} & \uline{\hfill Ю.В. Строганов \hfill} \\
			& \scriptsize (Подпись, дата) & \scriptsize (И.О. Фамилия)
		\end{signstabular}
		
		% \vspace{\baselineskip}
		
		% \begin{signstabular}[0.7]{p{7.25cm} >{\centering\arraybackslash}p{4cm} >{\centering\arraybackslash}p{4cm}}
			%     Консультант & \uline{\mbox{\hspace*{4cm}}} & \uline{\hfill  \hfill} \\
			%     & \scriptsize (Подпись, дата) & \scriptsize (И.О. Фамилия)
			% \end{signstabular}
	\end{table}
	
	\vfill
	
	\begin{center}
		\normalsize \textit{\textbf{2022} г.}
	\end{center}
\end{titlepage}
